%=====================================================================
\chapter{Reading and Critically Appraising a Paper}\label{ch:appraisal}
%=====================================================================
\locator{2}{Part~2 --- Finding and Synthesising Knowledge}

\begin{outcomes}
By the end of this chapter you will be able to:
\begin{tight}
  \item \textbf{Apply} the three-pass method and stop at the right pass for
    your purpose.
  \item \textbf{Interrogate} a paper's claim, evidence and limits rather than
    summarising it.
  \item \textbf{Take} notes that remain useful two years later.
  \item \textbf{Write} a referee report that is useful to authors and editor
    alike.
  \item \textbf{Reproduce} one claim from a paper and interpret what happens
    when you cannot.
\end{tight}
\end{outcomes}

\begin{prereq}
\chref{contribution} for the contribution types and the evidence each owes ---
appraisal is largely checking that match. \chref{slr} for where appraisal sits
in a review.
\end{prereq}

\begin{keyterms}
three-pass method $\bullet$ critical appraisal $\bullet$ claim extraction
$\bullet$ referee report $\bullet$ reproduction $\bullet$ annotated
bibliography $\bullet$ charitable reading
\end{keyterms}

\section{Reading Linearly Is the Problem}

The default way to read a paper --- start at the abstract, proceed to the
references, highlighting as you go --- is slow, and it produces a summary
rather than a judgement. Worse, it allocates the same effort to a paper you
will discard as to the one your thesis will build on. A graduate scholar who
must screen 900 records and deeply understand 30 of them cannot afford that.

The alternative is to read in passes of increasing depth, deciding after each
whether to continue.

\begin{paperbox}[The three-pass method]
S. Keshav, \emph{How to read a paper}, ACM SIGCOMM Computer Communication
Review 37(3), 2007~\cite{keshav2007read}.

\medskip
\textbf{Why it matters.} Two pages, and it changes how much literature you can
handle. It is the most immediately actionable thing in this part of the book.

\medskip
\textbf{What to extract.} The three passes and their time budgets, and his
description of using the method to do a literature survey --- which is a
different use of the same technique.
\end{paperbox}

\begin{center}
\small
\begin{longtable}{@{}L{1.5cm}L{2.0cm}L{4.6cm}L{5.6cm}@{}}
\toprule
\textbf{Pass} & \textbf{Budget} & \textbf{What you read} & \textbf{What you should be able to say} \\
\midrule
\endfirsthead
\toprule
\textbf{Pass} & \textbf{Budget} & \textbf{What you read} & \textbf{What you should be able to say} \\
\midrule
\endhead
\bottomrule
\endfoot
First  & 5--10 min
       & Title, abstract, introduction, section headings, conclusions,
         references (scanning for ones you know)
       & Category, context, correctness (plausible?), contributions, clarity.
         Enough to decide whether to read on \\
Second & About 1 hour
       & The whole paper, ignoring proofs and detailed derivations. Look
         carefully at every figure and table
       & The paper's claim and the evidence offered for it, in your own words.
         Enough to summarise it to somebody else \\
Third  & 4--5 hours
       & Everything, re-deriving the work: what would \emph{you} have done,
         and where do the authors' choices differ?
       & The paper's hidden assumptions, its unstated limitations, and what a
         reviewer should have objected to \\
\end{longtable}
\end{center}

\begin{conceptbox}[Which pass, for which purpose]
\begin{tight}
  \item \textbf{Screening for a review} (\chref{slr}): first pass only, on
    hundreds of papers. Anything more is a failure of discipline.
  \item \textbf{Building the related-work section:} second pass, on perhaps
    50. You need the claim and the evidence, not the derivation.
  \item \textbf{The three or four papers your own work stands on:} third pass,
    without exception. If your contribution is a delta from a paper, you must
    understand that paper better than its reviewers did.
  \item \textbf{Refereeing:} second pass minimum, third pass if you will
    recommend rejection. You owe an author whose work you are rejecting the
    effort of having understood it.
\end{tight}
\end{conceptbox}

\section{From Summarising to Interrogating}

A summary says what the paper claims. An appraisal says whether to believe it.
The difference is a set of questions asked in a fixed order.

\begin{center}
\small
\begin{tabular}{@{}L{0.6cm}L{13.8cm}@{}}
\toprule
1 & \textbf{What is the claim?} State it in one sentence, in your own words. If
    you cannot, that is a finding about the paper. \\
2 & \textbf{What kind of contribution is this?} (\chref{contribution}) The type
    determines what evidence is owed. \\
3 & \textbf{Is the evidence the right kind?} A causal claim from correlational
    data fails here, however elegant the analysis. \\
4 & \textbf{Is the evidence enough?} Sample size, power, number of
    datasets, number of runs (\chref{sampling}). \\
5 & \textbf{What is the comparison?} Against what baseline, and is it current
    and fairly tuned? \\
6 & \textbf{What is measured, and does it capture the construct?}
    (\chref{measurement}) \\
7 & \textbf{What would have falsified this?} If nothing, it was not a test. \\
8 & \textbf{Which threats do the authors name, and which do they not?} The
    unnamed ones are the interesting ones (\chref{validity}). \\
9 & \textbf{Could I reproduce this?} Are data, code and parameters available
    (\chref{reproducibility})? \\
10 & \textbf{What does this let \emph{me} do?} Method to borrow, baseline to
     beat, gap to fill, or nothing. \\
\bottomrule
\end{tabular}
\end{center}

\begin{alertbox}[Read charitably before you read critically]
There is a failure mode that looks like rigour: dismantling a paper for flaws
that are, on a fair reading, either acknowledged or irrelevant to its claim.

\medskip
The discipline is to first construct the strongest version of the paper's
argument --- including filling gaps the authors left implicit --- and only then
attack it. A criticism that survives that treatment is worth making. One that
does not was a misreading, and offering it in a seminar or a referee report
costs you credibility.

\medskip
This matters most when you disagree with a paper's conclusion. That is exactly
when your reading is least reliable.
\end{alertbox}

\begin{worked}[Second-pass appraisal of a plausible paper]
\textbf{The paper (constructed for this example).} ``Static Analysis Adoption
Reduces Post-Release Defects: Evidence from 200 Open-Source Projects.'' The
authors identify 200 GitHub projects, 100 that added a static analysis tool to
CI and 100 that did not, and report that adopters have 23\% fewer post-release
issues labelled \texttt{bug} in the following year ($p < 0.01$).

\medskip
\textbf{1--2. Claim and type.} Claim: adopting static analysis causes fewer
post-release defects. Type: empirical, and specifically causal.

\medskip
\textbf{3. Right kind of evidence?} No. This is observational: projects chose
whether to adopt. The claim is causal; the design supports association only.
This is the paper's central problem and everything else is secondary.

\medskip
\textbf{5. Comparison.} Non-adopters --- but adoption is not random. Teams that
add static analysis to CI are plausibly teams that already have CI, code
review, tests and maintainers who care. Any of those could produce the effect.
This is confounding by \emph{selection}, and \chref{quasiexp} exists to address
it: a difference-in-differences design using each project as its own control
before and after adoption would be far stronger.

\medskip
\textbf{6. Construct.} The outcome is issues labelled \texttt{bug}. Does that
measure defects, or does it measure \emph{labelling practice}? Teams
disciplined enough to adopt static analysis are plausibly also more disciplined
about triaging and labelling --- which could push the measured count in either
direction. The measure may be correlated with the treatment through a path that
has nothing to do with defects.

\medskip
\textbf{7. Falsifier.} As designed, almost nothing: any difference would have
been reported as an effect, and no difference would have been reported as a
null. The design does not discriminate between the causal hypothesis and the
selection hypothesis at all.

\medskip
\textbf{8. Threats named?} If the paper discusses selection bias, states the
labelling problem, and softens the claim to an association, it is a decent
descriptive study honestly reported. If it says ``reduces'' in the title and
recommends adoption in the conclusion, the claim outruns the evidence.

\medskip
\textbf{10. What it gives me.} A dataset construction method worth borrowing,
and a clearly identified gap: nobody has done this with an identification
strategy. That gap is a thesis.
\end{worked}

\begin{conceptbox}[Notice what the appraisal produced]
The worked example did not end in ``this paper is bad''. It ended in a research
opportunity that is specific, defensible, and directly derived from a named
weakness in existing work.

\medskip
That is what appraisal is for. \chref{contribution} said the delta is measured
against the literature; this is the mechanism by which you find it. A student
who appraises thirty papers this way finishes with a list of gaps, not a pile
of summaries.
\end{conceptbox}

\section{Notes That Survive}

\begin{pitfall}[Highlighting is not note-taking]
A highlighted PDF is unsearchable, tells you nothing about \emph{why} a passage
mattered, and is worthless in month 20 when you cannot remember which of 200
papers had the good threats section. Highlighting feels like work and produces
nothing.
\end{pitfall}

Write, for every paper you take past the first pass, a fixed short record ---
in your reference manager's note field, so it lives with the citation:

\begin{center}
\small
\begin{tabular}{@{}L{3.0cm}L{11.3cm}@{}}
\toprule
Claim          & One sentence, your words, not theirs \\
Type           & From \chref{contribution} \\
Method         & Design, subjects, N, measures \\
Result         & Effect size if given; ``not reported'' if not, which is itself
                 worth recording \\
Weakness       & The strongest objection you can make after a charitable
                 reading \\
Use to me      & Baseline / method / gap / background / discard \\
Quote          & At most one, verbatim, with a page number, for anything you
                 might quote later \\
RQ relevance   & Which of your research questions it bears on \\
\bottomrule
\end{tabular}
\end{center}

\begin{alertbox}[The one-quote rule protects you from yourself]
Record verbatim quotations \emph{as} quotations, in quotation marks, with a
page number, at the moment you take the note. The most common route to
inadvertent plagiarism is not dishonesty --- it is a paraphrase written into
notes six months ago that was actually a quotation, and is now indistinguishable
from your own prose. \chref{integrity} sets out what the consequences of that
look like under \reg{PhD Regs 2024}{\S32}, and they are severe enough that the
five seconds of discipline is worth it every time.
\end{alertbox}

\section{Refereeing, and Why You Should Volunteer}

Reviewing other people's work is the fastest way to improve your own, because
it puts you on the side of the table that decides. Ask your supervisor to
sub-review a paper for a conference. Most are glad to arrange it.

\begin{center}
\small
\begin{tabular}{@{}L{3.2cm}L{11.1cm}@{}}
\toprule
\textbf{Section} & \textbf{What goes in it} \\
\midrule
Summary        & The paper's claim and contribution, in your words. Authors
                 read this first to check you understood them --- and a
                 mistaken summary tells the editor to discount the rest \\
Strengths      & Genuine ones, specifically. Every paper has some \\
Major concerns & Numbered. Each states the problem, why it threatens the
                 claim, and what would resolve it \\
Minor points   & Typos, unclear figures, missing citations \\
Recommendation & With the reasoning, aimed at the editor \\
\bottomrule
\end{tabular}
\end{center}

\begin{conceptbox}[Two rules for a review you would be proud to sign]
\textbf{Every criticism carries an action.} ``The evaluation is weak'' is
useless. ``The evaluation compares against a 2016 baseline; the 2023 method
[ref] is stronger and available, and the claim of state-of-the-art performance
cannot stand without it'' is actionable and fair.

\medskip
\textbf{Write as though signed.} Whether or not the venue is open, write
nothing you would not put your name to. Anonymity is there to let you be
honest, not to let you be careless or unkind --- and the field is smaller than
it looks.
\end{conceptbox}

\section{The Deepest Reading Is Reproduction}

You do not fully understand an empirical or computational result until you have
tried to reproduce it. Take one claim --- not the whole paper --- and attempt
it from the artefacts provided.

\begin{center}
\small
\begin{tabular}{@{}L{4.6cm}L{9.7cm}@{}}
\toprule
\textbf{What happens} & \textbf{What it means} \\
\midrule
It reproduces
  & You now understand the method at a level no amount of reading gives you,
    and you have a working baseline for your own work \\
It reproduces after a struggle
  & Note exactly what was underspecified. That gap is publishable as a
    reproducibility study, and it is a real contribution
    (\chref{reproducibility}) \\
It does not reproduce
  & Before concluding anything: check your own setup twice, then write to the
    authors politely. Most non-reproduction is environment, version or seed,
    not error --- and authors are usually helpful \\
There are no artefacts
  & That is itself a finding about the field, and one worth counting in a
    review (\chref{slr}, quality item Q7) \\
\bottomrule
\end{tabular}
\end{center}

\begin{traditions}{appraising a paper}
\tradEMP{Check the design against the claim, the power against the effect, and
the measure against the construct. Ask what the confidence interval was, not
whether $p$ crossed a threshold.}
\tradDSR{Ask what the artefact demonstrates beyond its own existence, whether
the evaluation matches the claimed contribution, and whether the design
knowledge is stated separately from the implementation.}
\tradQUAL{Ask about the researcher's position, whether the description is thick
enough to interpret, whether interpretations are traceable to data, and whether
disconfirming cases were sought.}
\tradFORM{Check the proof. Then check the theorem statement --- most errors in
formal papers are in the assumptions, not the derivation, and an unrealistic
assumption makes a correct theorem irrelevant.}
\end{traditions}

\begin{lab}[Three passes and a referee report]
Choose a recent paper central to your topic.

\begin{tightnum}
  \item First pass, 10 minutes. Write the five C's: category, context,
    correctness, contributions, clarity.
  \item Second pass, one hour. Write the eight-field note record.
  \item Third pass, as long as it takes. Work through the ten appraisal
    questions and identify one assumption the authors did not state.
  \item Write a full referee report as if for the venue that published it,
    with a recommendation.
  \item Then attempt to reproduce one claim, and record what happened.
\end{tightnum}

\medskip
Bring the report to the week 5 seminar. Comparing reports on the same paper is
the fastest way to discover how differently competent readers weigh the same
evidence --- which is also why reviewers disagree.
\end{lab}

\begin{checkpoint}
\begin{tightnum}
  \item You have 900 records to screen and 30 papers to understand deeply. How
    many of each pass do you perform, and why is doing second passes during
    screening a mistake?
  \item Why should you construct the strongest version of a paper's argument
    before criticising it? Give a practical, not moral, reason.
  \item A paper reports $p = 0.03$ and no effect size. What do you write in the
    Result field of your note, and why does it matter later?
  \item You cannot reproduce a result. List, in order, the four things you do
    before concluding the paper is wrong.
\end{tightnum}
\end{checkpoint}

\begin{chaptersummary}
\begin{tight}
  \item Read in three passes and stop at the right one. Screening is first-pass
    work; the papers your delta is measured against get the third pass.
  \item Appraisal is ten questions, not a summary. The type of contribution
    determines what evidence is owed, and most weaknesses are a mismatch.
  \item Read charitably first. A criticism that survives the strongest version
    of the argument is worth making; one that does not was a misreading.
  \item Appraisal produces research opportunities, not verdicts. Thirty
    appraised papers give you a list of gaps.
  \item Keep a fixed eight-field note record in the reference manager.
    Highlighting produces nothing.
  \item Mark quotations as quotations, with page numbers, at the moment you
    write them. This is the main defence against inadvertent plagiarism.
  \item Referee something. Every criticism carries an action, and write as
    though signed.
  \item Reproduce one claim from the paper your work depends on. Whatever
    happens, you learn something worth having.
\end{tight}
\end{chaptersummary}

\begin{reviewq}
  \item The three-pass method allocates effort by expected value. Construct a
    case where it would fail --- a paper whose importance is invisible at first
    and second pass --- and say what would catch it.
  \item Peer review is criticised as slow, inconsistent and unable to detect
    fraud. Given that, why does the system persist, and what would a credible
    replacement have to provide?
  \item You referee a paper whose conclusion you believe is wrong but whose
    methods you cannot fault. What do you recommend, and what do you write?
  \item Reproduction is described here as the deepest form of reading, yet
    almost no graduate scholar does it. Identify the incentives that produce
    this, and propose one change to a course like this one that would shift
    them.
\end{reviewq}
